\section*{अध्याय \ref{ch:b3:partition-functions} --- सांख्यिकीय ऊष्मागतिकी: विभाजन फलन}

\begin{solution}{exo:b3:partition-functions:1}
हर वितरण को चारों दोलित्रों द्वारा रखे गए क्वांटमों के रूप में लिखिए, सबसे बड़ा
पहले। ``4, 0, 0, 0'': $4!/(3!\,1!) = 4$; ``3, 1, 0, 0'': $4!/(2!\,1!\,1!) = 12$; ``2, 2, 0, 0'':
$4!/(2!\,2!) = 6$; ``2, 1, 1, 0'': $4!/(1!\,2!\,1!) = 12$; ``1, 1, 1, 1'': 1। कुल 35 $= \binom73$।
सबसे प्रायिक ``3, 1, 0, 0'' और ``2, 1, 1, 0'' हैं (हर एक 12); दूसरा, जिसमें
0, 1, 2 क्वांटमों के लिए \omterm{def:b3:partition-functions:boltzmann}{समष्टियाँ} 1, 2, 1 हैं, पहले से ही ऊर्जा के साथ वैसे घटता है जैसे
बड़ी संख्याओं के लिए \omterm{def:b3:partition-functions:boltzmann}{बोल्ट्समान वितरण}।
\end{solution}

\begin{solution}{exo:b3:partition-functions:2}
$\eu^{-2500/(8.314 \times 298.15)} = \eu^{-1.009} = 0.365$; \qty{1000}{K} पर $\eu^{-0.301} = 0.740$। बहुत ऊँचे
ताप पर अनुपात 1 की ओर प्रवृत्त होता है (समान \omterm{def:b3:partition-functions:boltzmann}{समष्टियाँ}, कभी प्रतिलोमन नहीं)।
\end{solution}

\begin{solution}{exo:b3:partition-functions:3}
$m = 39.95\times10^{-3}/6.022\times10^{23} = \qty{6.634e-26}{kg}$; $\Lambda = h/\sqrt{2\pi mkT} =
\qty{1.600e-11}{m} = \qty{16.0}{pm}$; $q^{\mathrm T} = 10^{-3}/(1.600\times10^{-11})^3 = \num{2.44e29}$।
\end{solution}

\begin{solution}{exo:b3:partition-functions:4}
$hc/k = \qty{1.4388}{cm.K}$। \ce{N2}: $\theta_{\mathrm R} = \qty{2.863}{K}$, $q^{\mathrm R} = 298.15/(2 \times 2.863) = 52.1$।
\ce{HCl}: $\theta_{\mathrm R} = \qty{15.02}{K}$, $q^{\mathrm R} = 298.15/15.02 = 19.85$ (यथार्थ योग 20.19)।
\end{solution}

\begin{solution}{exo:b3:partition-functions:5}
\ce{I2}: $\theta_{\mathrm V} = 1.4388 \times 213.27 = \qty{306.9}{K}$, $q^{\mathrm V} = 1/(1 - \eu^{-1.029}) = 1.556$;
उत्तेजित अंश $1 - 1/q^{\mathrm V} = 0.357$। \ce{N2}: $\theta_{\mathrm V} = \qty{3352}{K}$, $q^{\mathrm V} = 1.000013$,
उत्तेजित अंश \num{1.3e-5}।
\end{solution}

\begin{solution}{exo:b3:partition-functions:6}
\ce{NO}: $q^{\mathrm E} = 2 + 2\eu^{-119.73/207.22} = 2 + 2 \times 0.561 = 3.12$; ऊपरी अंश $1.12/3.12
= 0.36$। उच्च ताप पर $q^{\mathrm E} \to 4$ (दोनों स्तर समान रूप से समष्टित)। \ce{Cl}:
$q^{\mathrm E} = 4 + 2\eu^{-882.35/207.22} = 4 + 2 \times 0.0142 = 4.03$।
\end{solution}

\begin{solution}{exo:b3:partition-functions:7}
$\ln q = \ln V + \frac32\ln T + \text{अचर}$, अतः $U - U(0) = NkT^2 \times \frac{3}{2T} = \frac32NkT$; एक
मोल के लिए \qty{298.15}{K} पर $\frac32RT = \qty{3.72}{kJ/mol}$, और $C_V = \frac32R = \qty{12.47}{J.K^{-1}.mol^{-1}}$।
\end{solution}

\begin{solution}{exo:b3:partition-functions:8}
$kT/p^\circ = \qty{4.116e-26}{m^3}$, $\Lambda^3 = \qty{4.093e-33}{m^3}$, अनुपात \num{1.006e7}:
$S_{\mathrm m} = 8.3145 \times (16.124 + 2.5) = \qty{154.85}{J.K^{-1}.mol^{-1}}$, सारणीबद्ध मान अंतिम
अंक तक। $p$ को 10 से भाग देने पर $R\ln10 = \qty{19.14}{J.K^{-1}.mol^{-1}}$ जुड़ता है।
\end{solution}

\begin{solution}{exo:b3:partition-functions:9}
$S^{\mathrm R}_{\mathrm m} = R[\ln(298.15/15.02) + 1] = 8.3145 \times (2.988 + 1) = \qty{33.16}{J.K^{-1}.mol^{-1}}$।
\end{solution}

\begin{solution}{exo:b3:partition-functions:10}
$\frac{\dd}{\dd J}\bigl[(2J+1)\eu^{-\theta J(J+1)/T}\bigr] = \bigl[2 - (2J+1)^2\theta/T\bigr]\eu^{-\theta J(J+1)/T} = 0$ से
$2J + 1 = \sqrt{2T/\theta}$, अर्थात् $J_{\max} = \sqrt{T/2\theta} - \frac12$। \ce{HCl}: 2.66, अतः $J = 3$; \ce{CO}:
6.86, अतः $J = 7$। P या \omterm{def:b3:rovibrational-spectroscopy:branches}{R शाखा} की किसी रेखा की तीव्रता उसके निचले स्तर की
\omterm{def:b3:partition-functions:boltzmann}{समष्टि} का अनुसरण करती है, जो $J_{\max}$ पर शिखर पर होती है, बैंड केंद्र पर स्थित रिक्ति से
कुछ रेखाएँ दूर।
\end{solution}

\begin{solution}{exo:b3:partition-functions:11}
$\int_0^\infty f\dd J = T/\theta$, $f(0) = 1$, $f'(0) = 2 - \theta/T$। अतः $q^{\mathrm R} \approx T/\theta + \frac12 -
\frac16 + \frac{\theta}{12T} \approx T/\theta + \frac13$। $T/\theta$ की आपेक्षिक त्रुटि लगभग $\theta/3T$ है:
1\,\% से कम जब $T > 33\,\theta$। \ce{H2} के लिए $\theta = 1.4388 \times 59.32 = \qty{85.3}{K}$:
\qty{298.15}{K} पर त्रुटि 10\,\% है, स्वीकार्य नहीं (और $J$ पर नाभिकीय-चक्रण
प्रतिबंध को भी यथार्थ रूप से लेना होगा, \cref{ch:b3:statistical-thermo-applied})।
\end{solution}

\begin{solution}{exo:b3:partition-functions:12}
$x = \theta_{\mathrm V}/T \to \infty$ के लिए $x/(\eu^x - 1) \to 0$ और $\ln(1 - \eu^{-x}) \to 0$: $S \to 0$। $x \to 0$ के लिए
$x/(\eu^x - 1) \approx 1 - x/2$ और $-\ln(1 - \eu^{-x}) \approx -\ln x + x/2$: $S \to R(1 - \ln x) = R[1 + \ln(T/\theta_{\mathrm V})]$।
\ce{I2}, $x = 1.029$: यथार्थ $R(0.5721 + 0.4420) = \qty{8.43}{J.K^{-1}.mol^{-1}}$; सीमा $R(1 - 0.0289) =
\qty{8.07}{J.K^{-1}.mol^{-1}}$। $T \approx \theta_{\mathrm V}$ पर भी चिरसम्मत सीमा केवल 4\,\% कम है।
\end{solution}

\begin{solution}{pb:b3:partition-functions:1}
\textbf{1.} $m = 28.014\times10^{-3}/6.02214\times10^{23} = \qty{4.652e-26}{kg}$।
\textbf{2.} $\Lambda = h/\sqrt{2\pi mkT} = \qty{1.910e-11}{m}$ (\qty{19.1}{pm})।
\textbf{3.} $kT/p^\circ = 1.380649\times10^{-23} \times 298.15/10^5 = \qty{4.116e-26}{m^3}$।
\textbf{4.} $q^{\mathrm T}/N = 4.116\times10^{-26}/(1.910\times10^{-11})^3 = \num{5.905e6}$: प्रति अणु लगभग साठ लाख
स्थानांतरणीय अवस्थाएँ, अतः दो अणु लगभग कभी एक अवस्था साझा नहीं करते और
$Q = q^N/N!$ मान्य है।
\textbf{5.} $S^{\mathrm T}_{\mathrm m} = R(\ln\num{5.905e6} + \frac52) = 8.3145 \times 18.091 = \qty{150.42}{J.K^{-1}.mol^{-1}}$।
\textbf{6.} $+R\ln2 = \qty{5.76}{J.K^{-1}.mol^{-1}}$, अर्थात् \num{156.18}।
\textbf{7.} $B_0 = 1.998241 - 0.008659 = \qty{1.989582}{cm^{-1}}$।
\textbf{8.} $\theta_{\mathrm R} = 1.438777 \times 1.989582 = \qty{2.8626}{K}$।
\textbf{9.} $\sigma = 2$: दोनों नाभिक सर्वसम हैं; अणु को सिरे से सिरे तक उलटने पर
अविभेद्य अभिविन्यास मिलता है, और नाभिकों की विनिमय सममिति,
हर नाभिकीय-चक्रण अवस्था के लिए, केवल आधे घूर्णी स्तरों की अनुमति देती है।
\textbf{10.} $q^{\mathrm R} = 298.15/(2 \times 2.8626) = 52.08$।
\textbf{11.} हाँ: $T/\theta_{\mathrm R} = 104$, $q^{\mathrm R}$ पर लगभग 0.3\,\% की त्रुटि और
एन्ट्रॉपी पर उससे बहुत कम।
\textbf{12.} $S^{\mathrm R}_{\mathrm m} = R(\ln52.08 + 1) = 8.3145 \times 4.9527 = \qty{41.18}{J.K^{-1}.mol^{-1}}$।
\textbf{13.} $J_{\max} = \sqrt{298.15/5.725} - 0.5 = 6.72$: $J = 7$ (\omterm{def:b3:partition-functions:boltzmann}{समष्टि} 0.0839, $J = 6$ के
0.0831 से थोड़ी ऊपर)।
\textbf{14.} $2358.57 - 2 \times 14.324 = \qty{2329.92}{cm^{-1}}$।
\textbf{15.} $\theta_{\mathrm V} = 1.438777 \times 2329.92 = \qty{3352}{K}$।
\textbf{16.} $x = 3352.2/298.15 = 11.24$; $q^{\mathrm V} = 1/(1 - \eu^{-11.24}) = 1.000013$।
\textbf{17.} $\eu^{-11.24}/q^{\mathrm V} = \num{1.3e-5}$।
\textbf{18.} $S^{\mathrm V}_{\mathrm m} = R[11.24\eu^{-11.24} + \eu^{-11.24}] = \qty{0.0013}{J.K^{-1}.mol^{-1}}$, नगण्य।
\textbf{19.} $g_0 = 1$ (${}^1\Sigma$): $R\ln1 = 0$। उत्तेजित इलेक्ट्रॉनिक अवस्थाएँ कई
इलेक्ट्रॉनवोल्ट ऊपर हैं, $kT$ के सौ गुने से अधिक: उनके बोल्ट्समान गुणक
पूर्णतः नगण्य हैं।
\textbf{20.} $150.42 + 41.18 + 0.00 + 0 = \qty{191.60}{J.K^{-1}.mol^{-1}}$।
\textbf{21.} अंतर \qty{0.01}{J.K^{-1}.mol^{-1}} है, किए गए
सन्निकटनों (\omterm{def:b3:quantum-model-systems:rotor}{दृढ़ घूर्णक}, गतियों का पृथक्करण) के स्तर पर।
\textbf{22.} कुछ केल्विन से मापी गई ठोस की ऊष्मा धारिताएँ (उसके नीचे $T^3$ के रूप में
बहिर्वेशित), ठोस--ठोस संक्रमण, गलन और
वाष्पन की एन्थैल्पियाँ उनके तापों से विभाजित, गैस की ऊष्मा धारिता,
और अनादर्शता के लिए एक संशोधन।
\textbf{23.} नाभिक हर अभिक्रिया में संरक्षित रहते हैं, अतः उनकी चक्रण एन्ट्रॉपी
दोनों ओर समान है और निरस्त हो जाती है; ऊष्मामिति भी उसे नहीं देखती, क्योंकि
नाभिकीय चक्रण क्रिस्टल में मापे गए सबसे निम्न तापों तक अव्यवस्थित
बने रहते हैं।
\textbf{24.} \textbf{स्पेक्ट्रमी स्थिरांकों से $S^\circ(\ce{N2}, \qty{298.15}{K}) = \qty{191.6}{J.K^{-1}.mol^{-1}}$},
जिसमें 150.4 स्थानांतरण से और 41.2 घूर्णन से।
\end{solution}
